% !TeX program = lualatex
% Compile twice: lualatex 169.tex
% All references and prose are embedded; no BibTeX database or external inputs.
\documentclass[11pt,a4paper]{article}
\usepackage[margin=24mm,headheight=14pt]{geometry}
\usepackage{fontspec}
\setmainfont{Latin Modern Roman}
\setsansfont{Latin Modern Sans}
\setmonofont{Latin Modern Mono}
\usepackage{amsmath,amssymb,mathtools}
\usepackage{unicode-math}
\setmathfont{Latin Modern Math}
\usepackage{microtype}
\usepackage{enumitem,xurl,needspace}
\usepackage[dvipsnames]{xcolor}
\usepackage{hyperref}
\hypersetup{unicode=true,colorlinks=true,linkcolor=MidnightBlue,urlcolor=MidnightBlue,
 pdftitle={500 Mathematical Problem Statements},pdfauthor={Source-annotated compilation},bookmarksnumbered=true}
\usepackage{bookmark}
\usepackage{tocloft}
\setlength{\cftsecnumwidth}{3em}
\cftsetpnumwidth{2.5em}
\cftsetrmarg{4em}
\usepackage{titlesec}
\titleformat{\section}{\Large\bfseries\raggedright}{\thesection}{0.7em}{}
\usepackage{fancyhdr}
\pagestyle{fancy}\fancyhf{}
\fancyhead[L]{\small\sffamily Mathematical problem statements}
\fancyhead[R]{\small Source edition 22}
\fancyfoot[C]{\thepage}
\setlength{\parindent}{0pt}
\setlength{\parskip}{5pt plus 1pt}
\setlength{\emergencystretch}{3em}
\setcounter{tocdepth}{1}
\setcounter{secnumdepth}{1}
\setlist[itemize]{leftmargin=3em,itemsep=5pt,topsep=4pt,parsep=0pt}
\allowdisplaybreaks
\newcommand{\N}{\mathbb{N}}
\newcommand{\Z}{\mathbb{Z}}
\newcommand{\Q}{\mathbb{Q}}
\newcommand{\R}{\mathbb{R}}
\newcommand{\C}{\mathbb{C}}
\newcommand{\F}{\mathbb{F}}
\newcommand{\A}{\mathbb{A}}
\newcommand{\PP}{\mathbb P}
\newcommand{\E}{\mathbb{E}}
\newcommand{\eps}{\varepsilon}
\newcommand{\dd}{\,\mathrm d}
\newcommand{\sref}[2]{\hyperlink{source:#1:#2}{[S#2]}}
\newcommand{\eref}[2]{\hyperlink{extra:#1:#2}{[E#2]}}
% Turn the next switch off for a shorter reading copy. The quotations preserve
% every exactTarget field, including qualifications omitted from a short summary.
\newif\ifshowcatalogue
\showcataloguetrue
\newcommand{\cataloguescope}[1]{\ifshowcatalogue
 \paragraph{Exact scope in the supplied catalogue.}
 \begin{quote}\small #1\end{quote}\fi}

% Standalone notebook wrapper; original problem section retained verbatim.
\numberwithin{equation}{section}
\hypersetup{pdftitle={Research attempt -- problem 169},
 pdfauthor={Research notebook; source compilation retained}}
\fancyhead[L]{\small\sffamily Research notebook}
\fancyhead[R]{\small\sffamily Problem 169 | 27 September 2026}
\begin{document}
\setcounter{section}{168}
% ================================================================
% problemId: problem.weinstein-conjecture-on-periodic-orbits-of-reeb-flows
% source releaseRank: 169; source releaseStatus: open_with_solved_subcases
% exactTarget SHA-256: 5fbbd3de181dbfc77a628c8ade873cb42a0e006dc895ba3cbe5eb043e4c5b086
\clearpage
\section{\texorpdfstring{Weinstein Conjecture on periodic orbits of Reeb flows}{Weinstein Conjecture on periodic orbits of Reeb flows}}\label{169}
\subsection{Definitions and mathematical statement}
On a closed smooth $(2n+1)$-manifold $M$, a contact form is a one-form $\alpha$ with $\alpha\wedge(\dd\alpha)^n$ nowhere zero. Its Reeb vector field $R_\alpha$ is uniquely specified by
\[
 \alpha(R_\alpha)=1,\qquad \iota_{R_\alpha}\dd\alpha=0.
\]
A periodic orbit is a smooth map $\gamma:\R/T\Z\to M$, $T>0$, satisfying $\dot\gamma=R_\alpha(\gamma)$. The conjecture asserts the existence of such an orbit for every closed contact manifold and every contact form in all dimensions. Here closed means compact without boundary. No nondegeneracy, contractibility, or uniform bound on the period is imposed.
\par\smallskip\textit{Formulation and concept references:} \sref{169}{1}.
\cataloguescope{Prove that every closed contact manifold admits a periodic orbit of the Reeb flow of any contact form, in all dimensions.}
\subsection{Short English statement}
Must every Reeb flow on a compact contact manifold without boundary have at least one closed trajectory?
% BEGIN RESEARCH ADDITION ATTEMPT-169
\subsection{Research attempt}

\paragraph{Current frontier.}
The smooth closed-manifold assertion in the question is the one recorded in
\sref{169}{1}, Section 5.1; results for singular contact forms in that paper
must not be substituted for it. Taubes's theorem settles every contact form
in dimension three. Hutchings gives both its precise statement and a survey
of higher-dimensional special cases \eref{169}{1}, Theorem 1.3 and
Sections 2--3. Those cases do not include every closed contact manifold in
all dimensions. The investigation below does not assume nondegeneracy,
fillability, a contractible orbit, or a bound on the period in the target.

\paragraph{Attack route.}
A holomorphic-curve approach requires a nonvanishing invariant and compactness
in the relevant contact geometry; a return-map approach needs a suitable
section and a fixed-point mechanism. A third possibility is to approximate
an arbitrary form by forms with explicit periodic fibers and pass to a
limit. The decisive issue for that approach is a locally uniform upper
bound on the selected periods. We prove the needed compactness lemma,
construct controlled orbits for a restricted but infinite-dimensional class,
and test whether either ingredient extends to arbitrary perturbations.

\paragraph{Attempt: exactly what approximation would need.}
Fix an auxiliary Riemannian metric on the compact manifold. Suppose that
contact forms $\alpha_j$ converge to a contact form $\alpha$ in $C^2$.
Solving the pointwise nonsingular linear equations defining the Reeb field
shows that $R_{\alpha_j}\to R_\alpha$ in $C^1$. Suppose further that there
are periodic trajectories with initial points $p_j$ and periods
$0<T_j\le C$, where $C$ is independent of $j$. Then $\alpha$ has a
periodic trajectory.

Here are the two compactness details that this statement needs. Because
$R_\alpha$ is nowhere zero and the manifold is compact, its norm has a
positive minimum. In a finite collection of coordinate charts, the nearby
fields have a common positive lower speed and common bounds on their values
and first derivatives. The integral equation for their flows gives,
uniformly for sufficiently small $t$,
\[
 \phi_j^t(p)-p=tR_{\alpha_j}(p)+O(t^2).
\]
Choose $t_0>0$ so that the remainder cannot cancel the first term and the
flow remains in the chosen chart. Thus there are no periods in $(0,t_0)$,
uniformly in large $j$. Passing to a subsequence gives $p_j\to p$ and
$T_j\to T\in[t_0,C]$. Continuous dependence of flows on the vector field,
uniform on bounded time intervals, implies
\[
                    \phi_\alpha^T(p)=p.
\]
The resulting trajectory is nonconstant since its velocity never vanishes.
This proves the lemma, including the lower-period issue.

Consequently a sufficient approximation statement is: for every $\alpha$
there is a $C^2$-convergent sequence of contact forms with periodic orbits
whose periods are bounded by a constant depending on $\alpha$. Density of
forms possessing an orbit, with no control of its period, does not establish
this sufficient statement. The divergence $T_j\to\infty$ is a genuine
failure of this compactness argument, even though the trajectories stay in
a compact manifold.

\paragraph{Attempt: controlled periodic fibers under basic rescaling.}
Consider a regular contact form $\alpha_0$ whose Reeb flow is a free circle
action with period $T_0$. Write $\pi:M\to B$ for the quotient and
$d\alpha_0=\pi^*\omega$, where $\omega$ is symplectic. This is the regular
circle-bundle setting discussed in \eref{169}{1}, Section 2. Let $f$ be
\emph{any} smooth positive function on $B$, and put
$\beta=(f\circ\pi)\alpha_0$. Suppress pullbacks in the following calculation.
There is a unique horizontal vector field $Y$ such that
\begin{equation}\label{eq169-basic}
 R_\beta=f^{-1}R_{\alpha_0}+Y,
 \qquad \iota_Yd\alpha_0\big|_{\ker\alpha_0}=f^{-2}df.
\end{equation}
Indeed the first formula gives $\beta(R_\beta)=1$. For a horizontal $V$,
\[
 d\beta(R_\beta,V)=-f^{-1}df(V)+f\,d\alpha_0(Y,V)=0.
\]
The remaining direction vanishes as well: evaluating the second equation
on $Y$ gives $df(Y)=0$, while $df(R_{\alpha_0})=0$. Nondegeneracy on the
horizontal distribution proves uniqueness and the Reeb equations.

At a critical point $b$ of $f$, the field $Y$ vanishes on the entire fiber.
That fiber is therefore periodic, with period $T_0 f(b)$. Taking a minimum
of $f$ gives the explicit bound
\begin{equation}\label{eq169-bound}
                     T\le T_0\min_B f.
\end{equation}
The calculation is not perturbative and requires no Morse hypothesis on
$f$. It also checks the normalization under constant rescaling, for which
Reeb time is multiplied by the same constant. Thus the desired compactness
bound is available in this restricted class.

The attempt to extend the construction fails at a precise point. For a
nonbasic function $h$ on $M$, a critical point of $h$ is not an entire
critical circle; the horizontal part of $R_{h\alpha_0}$ need not vanish
along its original fiber. Averaging $h$ along the circle action produces
a basic function but is not an arbitrarily small approximation operation.
The basic functions form a closed proper subspace in $C^2$: every such
function satisfies $R_{\alpha_0}h=0$, and this equation survives limits.
A function with a nonzero derivative in the fiber direction cannot be
approximated in $C^2$ by basic ones. No density assertion follows from
\eqref{eq169-basic}.

\paragraph{Attempt: testing a global return section.}
There is also an elementary obstruction to a tempting replacement argument.
For $n\ge1$, a nonempty closed hypersurface $\Sigma^{2n}\subset M$ cannot
be everywhere transverse to the Reeb field. If it were, $d\alpha|_\Sigma$
would be symplectic: the only kernel of $d\alpha$ in $TM$ is the Reeb line,
and transversality identifies $T\Sigma$ with the symplectic quotient by
that line. With its symplectic orientation,
\[
 0<\int_\Sigma(d\alpha)^n
   =\int_\Sigma d\bigl(\alpha\wedge(d\alpha)^{n-1}\bigr)=0,
\]
a contradiction. Thus an everywhere-transverse compact return section
without boundary is unavailable. Sections with boundary on periodic orbits
are not ruled out, but assuming such a boundary would already assume the
existence sought. This is an obstruction to that proof device, not to the
Weinstein conjecture.

\paragraph{Stress test.}
The dimension-one case is separate and immediate: each component of a closed
one-manifold is a circle, and a nowhere-zero Reeb field traverses it in a
finite positive time. The Stokes argument starts only in dimension three.
The controlled-fiber construction needs a regular contact form to start
with; neither existence of that form for a given contact structure nor
approximation of arbitrary contact structures by regular ones has been
proved here. Likewise, a general volume-preserving aperiodic flow is not a
counterexample unless it satisfies the contact Reeb equations.

In the limit lemma the upper-period bound is an additional hypothesis, not
a hidden conclusion of compactness of $M$. Allowing $\alpha_j$ to converge
to a noncontact form would also invalidate the nonsingular linear-system
argument. Every assumption used in that lemma is therefore kept explicit.
The geometric computations are deductions within this attempt using a
standard circle-bundle example; no claim of a newly solved higher-dimensional
case or of literature priority is made.

\paragraph{Outcome.}
\textbf{Outcome: Exploratory approach with unresolved gap.}

The attempt proves a bounded-period closure lemma, an exact period formula
for positive basic rescalings of a regular contact form, and an obstruction
to a closed global-section shortcut. These identify a possible approximation
route and show why its easiest implementation does not reach arbitrary
forms. A general construction of approximants with uniformly bounded
periods, or a different existence mechanism covering the omitted contact
geometries, remains unproved.
% END RESEARCH ADDITION ATTEMPT-169
\subsection{Sources}
\begin{itemize}\raggedright
\item[\textnormal{[S1]}]\hypertarget{source:169:1}{} Eva Miranda and Cedric Oms, The singular Weinstein conjecture, Advances in Mathematics 389 (2021), Section 5.1, Conjecture 5.1.
\url{https://arxiv.org/pdf/2005.09568}.
{\small\textit{Catalogue formulation source.}}
% BEGIN RESEARCH ADDITION SOURCES-169
\item[\textnormal{[E1]}]\hypertarget{extra:169:1}{} Michael Hutchings,
\emph{Taubes's proof of the Weinstein conjecture in dimension three},
author survey, Theorem 1.3 and Sections 2--3.
\url{https://math.berkeley.edu/~hutching/pub/tw.pdf}.
{\small\textit{Consulted 27 September 2026. Used for the three-dimensional
theorem and regular contact circle-bundle examples, not an all-dimensional
existence theorem.}}
% END RESEARCH ADDITION SOURCES-169
\end{itemize}

\end{document}
